\chapter{Introdução}
\label{cap:introducao}

% Os algoritmos de aprendizado de máquina que fazem classificação em imagens, necessitam que estas sejam de baixa dimensionalidade \cite{yang2019re} . \gls{KNN} askjdhkjasdhj ashdkajsdhajkhj \gls{KNN}

No cenário contemporâneo de digitalização dos serviços, o setor médico-veterinário enfrenta o desafio de transpor a eficiência clínica para o ambiente domiciliar. Essa transição é impulsionada não apenas pela demanda por conveniência, mas sobretudo pelo bem-estar animal. Ao analisar o comportamento felino, observou-se que "os gatos apresentaram níveis de estresse significativamente menores quando examinados em um ambiente doméstico em comparação com um ambiente clínico", (Nibblett, 2015). Diante disso, torna-se imperativo o desenvolvimento de ferramentas robustas voltadas especificamente para a gestão do atendimento remoto e volante.

Tendo isso em vista, médicos veterinários que atuam em atendimentos domiciliares enfrentam desafios acentuados que comprometem a agilidade e a segurança do serviço. A dependência de conectividade constante cria uma barreira crítica: em localidades com infraestrutura de rede deficitária, o uso de sistemas baseados exclusivamente na presença de internet e nuvem, frequentemente inviabiliza o fluxo ideal, forçando o veterinário a recorrer a registros improvisados em papel ou na memória. Estudos na área de informática em saúde apontam que essa fragmentação da informação eleva significativamente o risco de perda de dados, problemas de legibilidade e erros de transcrição posterior (F.L Pollari, 1996). Além das falhas operacionais, essa dinâmica impede o acesso imediato ao histórico clínico completo do paciente no local de atendimento, momento crucial do exame clínico no qual a tomada de decisão exige rapidez e fundamentação baseada em evidências (Kristine M. Alpi, 2011).

De que foram uma arquitetura \textit{offline-first} pode contribuir para a sincronização de dados clínico-veterinários em contextos de baixa conectividade? Diante desse problema, é necessária a investigação de técnicas que contribuam com o armazenamento de dados de uma aplicação veterinária, de forma que, sem a presença de internet no dispositivo móvel do usuário, os dados sejam persistidos primeiramente na camada local e, posteriormente, no servidor. Para alcançar esse objetivo, o seguinte trabalho apresenta o desenvolvimento do RoadVet, uma aplicação móvel voltada para médicos veterinários autônomos que atuam em atendimento domiciliar, baseada na arquitetura \textit{offline-first}, arquitetura que prioriza o uso da aplicação no dispositivo do usuário, buscando que os dados sejam salvos primordialmente na camada de apresentação, que no presente trabalho é implementada com a tecnologia \textit{React Native} para a interface do usuário e a biblioteca \textit{WatermelonDB}, responsável por persistir esses dados no dispositivo do veterinário, quando a conexão com a internet é identificada, os dados que foram criados são sincronizados com a camada de suporte, que foi abstraída como uma API(\textit{Application Programming Interface}), construída utilizando o\textit{ Django Rest Framework} e um banco de dados relacional\textit{ SQLite}, para manter as informações acessíveis.

Além desta introdução, o presente documento está estruturado da seguinte maneira: o Capítulo ~\ref{cap:fundamentacao-teorica} trata da fundamentação teórica, abordando o mercado de cuidado animal e o meio social, o paradigma \textit{offline-first} e a legislação que apoia esses animais de estimação. O capítulo ~\ref{cap:trabalhos-relacionados} traz os trabalhos relacionados que serviram como exemplo para o desenvolvimento do RoadVet. O capítulo ~\ref{chap:metodologia} destaca a metodologia adotada, que envolve a caracterização da pesquisa, discussão com profissionais da área veterinária, processo de levantamento de requisitos, com alguns exemplos, definição da arquitetura da aplicação que foi desenvolvida e o procedimento de validação e testes adotados. O capítulo ~\ref{cap:roadvet} trata do detalhamento do desenvolvimento do RoadVet, demonstrando visualmente o fluxo da aplicação, entre\textit{ back-end} e \textit{front-end}, uma visão geral do aplicativo móvel e a especificação do fluxo de uso principal do RoadVet. O capítulo ~\ref{cap:conclusao} debate as conclusões, dificuldades encontradas e os trabalhos futuros.

% \ref{fig:fluxogramaDeAlgumaCoisa}. 

% \ref{fig:oficial}

\section{Motivação}
\label{sec:motivacao}

Neste contexto, percebe-se que o uso de ferramentas tecnológicas que não dependam de acesso a Internet podem ajudar o trabalho de um veterinário domiciliar. Ao desenvolver uma solução que funcione de maneira \textit{offline}, busca-se otimizar o tempo desses profissionais, ao mesmo tempo que oferecer um serviço que aumente a eficiência e qualidade do cuidado animal, permitindo que o foco permaneça no atendimento e não nas limitações tecnológicas do local.

Este desenvolvimento estabelece, ainda, o alicerce necessário para futuras expansões modulares, como a otimização logística de deslocamentos.

\section{Objetivos}
\label{sec:objetivos}

Este trabalho tem como objetivo abordar o tema proposto e identificar as principais necessidades de médicos veterinários. A seguir, são apresentados os objetivos geral e específicos que norteiam este trabalho.

\subsection{Objetivo Geral}
\label{sec:objetivo-geral}

Desenvolver uma aplicação móvel voltada para atendimentos médico-veterinários domiciliares, utilizando uma arquitetura \textit{offline-first} capaz de persistir localmente e sincronizar no servidor, os dados, mesmo em ambientes de baixa conexão com a internet.

\subsection{Objetivos Específicos}
\label{sec:objetivos-especificos}

Os objetivos específicos do presente trabalho consistem na implementação da estratégia \textit{offline-first}, que foca na persistência de dados, partindo da camada do cliente. Essa abordagem visa garantir que os dados inseridos pelo usuário sejam sincronizados no servidor. Para tanto definem-se os seguintes passos: 

	\begin{itemize}
     \renewcommand\labelitemi{-}
		\item Analisar e implementar tecnologias existentes que sejam capazes de persistir os dados localmente, em dispositivos móveis.
		\item Implementar tecnologia de armazenamento local na camada do usuário, de forma que sejam visíveis os dados que ainda não foram sincronizados com o servidor.
        \item Projetar mecanismo de sincronização assíncrona.
		\item Realizar simulação de testes clínicos reais, de forma que o usuário final tenha uma arquitetura segura.
	\end{itemize}